\chapter{Who Speaks, and Who Sees?}
\section{A Written Self-Criticism}
Pick up a sheet of paper.

A composition sheet with space for four hundred Chinese characters bears ``Self-Criticism'' neatly centred at the top, the handwriting almost excessively proper. It begins:

``With profound sorrow, I reflect on my misconduct on the basketball court yesterday afternoon. I deeply recognise that my recklessness harmed a fellow student and disgraced our class\ldots''

Its writer, Xiaoming, is in eighth grade, hunched at his desk biting his pen. Look closely and you notice something awkward. His face shows no profound sorrow, only impatience. At ``disgraced our class'', he pauses and mutters, ``Is it really that serious?''

Ten minutes earlier, he messaged a close friend in quite different language: ``I only went to save the ball. That younger kid lost his balance, and I didn't step on his glasses. Why should I write this? Besides, he said he was fine.''

Hold the paper farther away and consider something usually overlooked: who is its ``I''?

The real, pen-biting Xiaoming, full of objections? Apparently not. He thinks he did nothing wrong, while the paper's ``I'' deeply recognises wrongdoing. The Xiaoming who collided with someone on court? Not quite either. That person appears once, pinned beneath ``recklessness''. The half-second hesitation while saving the ball and his first words while helping the boy up appear nowhere.

Three people inhabit this paper: the writer, the speaking ``I'', and the colliding ``I'' described by that voice. They share a name but are not the same person.

This is not Xiaoming's cunning or hypocrisy peculiar to self-criticism. It is a structure within every story and account, unusually exposed here. We will dismantle it: when a story is told, who speaks, who sees, and who has not been permitted to speak?

\section{Three Seconds on the Court}
Return to the scene and slow three seconds down.

Thursday afternoon, after third period, on the outdoor court east of the school. Late autumn's low sun stretches the basket's shadow. Yesterday's rain has left a small wet patch on the concrete. Six boys play half-court. Uniform jackets hang from the wire fence; two seventh-grade classes have PE nearby, whistles sounding intermittently.

Three against three, a close score. An opponent throws a pass too hard. The ball flies out towards a seventh-grade spectator at the sideline, a small boy with glasses holding a stack of newly distributed exercise books.

Xiaoming, nearest the ball, turns and chases. In slow motion, he takes two long strides, reaches it with his fingertips, and knocks it back into play. But he cannot arrest his momentum. His shoulder hits the boy squarely. The boy falls backwards into a sitting position; books scatter, glasses fly half a metre, and concrete scrapes his knee.

Three seconds altogether.

The next twenty minutes unfold not on the court but from mouth to mouth.

The injured boy is helped towards the infirmary. On the way, he tells a companion, ``He just came straight at me, as if he hadn't seen me.''

At the sideline, Xiaoming tells the gathering classmates, ``I reached the ball. It was momentum; nobody could have stopped.''

A girl from the neighbouring class tells her desk-mate, ``It was a hard hit, but I don't think the older boy meant it. He helped pick up the books.''

The sports representative runs to the class teacher. His opening is, ``Teacher, Xiaoming knocked over a seventh-grader.'' Notice what disappears: hesitation, momentum, the wet patch, the flying ball. Only one person, one action, and one victim remain.

By evening, the story reaches a year-group chat whose members were absent: ``Have you heard? An eighth-grader knocked a seventh-grader into the infirmary while playing basketball.''

Nobody lies. Almost every sentence has a corresponding fact. Yet one collision in one second grows into five stories that do not fit together.

Before judging which is right, ask an earlier question: through whose eyes did each story grow?

\section[Three Selves and a Moving Camera]{Three Selves and a Camera That Moves}
Now lay out the conflict.

The class teacher has reasons for requiring self-criticism; Xiaoming has reasons against it. The teacher says writing makes reflection heartfelt. Xiaoming thinks every written word is false while he remains unconvinced, so what is such reflection worth?

Both sound reasonable. But their disagreement is not really whether a collision deserves an apology. It is something deeper: \textbf{when someone tells their own story on paper, who is the paper's ``I''?}

Narratology, the study of how stories are told, supplies three terms we must separate. Do not fear them; the self-criticism makes each clear.

\textbf{Author}: the real person writing the words. Pen-biting Xiaoming is the author, flesh and blood, impatient, sending complaints by message.

\textbf{Narrator}: the voice speaking inside the text. The profoundly sorrowful ``I'' who deeply recognises error is the narrator. He exists only on paper; words build his emotions and sincerity. The craft of this assignment is for real Xiaoming to construct a properly contrite ``I'' to hand to the teacher.

\textbf{Character}: the agent described in the story. The Xiaoming who recklessly collided with someone yesterday is a character, an image assembled through the narrator's selection. The half-second saving the ball is removed; the half-second collision remains.

Three people, one name. Xiaoming's grievance now fits a sentence: the author disagrees with the narrator, who tells only half the character. Complaining in messages, he speaks as author; writing self-criticism, he plays narrator. When others say he knocked someone over, he becomes a character in their story, with no right even to speak, only to be struck, described, and classified.

Author, narrator, and character are not the same person. This is not a literary game's rule but the first key to understanding every account, including those not beginning with ``I''.

Narrative person concerns the narrator's distance from the story. Three common forms are:

\textbf{First person}: an ``I'' speaks while living within the story. Self-criticisms, diaries, and voice messages belong here. Its strength is closeness; its limit is what this ``I'' can see.

\textbf{Third person}: the teller stands outside, calling everyone by name, Xiaoming or the seventh-grade boy. Most histories and news reports use this form. It appears objective, but its outside teller still chooses where the camera points and which half-second disappears.

A more slippery form is \textbf{free indirect discourse}. Consider:

``Xiaoming stood outside the office. Now he was done for. Teacher Wang must know everything.''

The first sentence is standard third-person narration. Whose voice says ``done for''? Not the outside narrator's, but Xiaoming's inward voice. Without quotation marks or ``he thought'', it enters narration, overlapping the narrator's voice with the character's thoughts. You seem to read someone else's account while quietly entering a character's mind. Novels abound in this covert movement. Watch for it next time; it may surprise you.

How old is this distinction? Consider China's most famous novel. \emph{Dream of the Red Chamber} opens with a small poem:

\begin{quote}
Pages full of absurd words, a handful of bitter tears. All call the author foolish; who understands their flavour? (Chapter 1)
\end{quote}
More intriguingly, the novel claims the story was originally carved on a great stone, copied by a Daoist, and later read for ten years and revised five times by Cao Xueqin in his Mourning-the-Red Studio. Its narrator thus claims Cao is not the author but an editor. The real Cao Xueqin is, of course, the author, yet he makes a stone the first narrator and demotes himself to copyist. Three hundred years ago, writers already played the author-is-not-narrator mechanism, more deeply than we do.

\section{Everyone Has a Case}
Return to the court and present each position fully before judging.

\textbf{The injured boy's case}: I stood outside the boundary and bothered nobody. No one warned me before the ball arrived. Saving it was your choice, but my knee and glasses paid. Afterwards everyone gathered around the players debating intention, as though I were a prop. I do not want slow-motion analysis of three seconds. I want someone first to ask how I am.

\textbf{The spectators' case}: we stood farthest away and saw most. He reached the ball, yes, but never glanced at the sideline before starting. Not intending harm and having no responsibility differ; drivers do not intend rear-end collisions either.

\textbf{Now give Xiaoming's full case}. This deserves care because by the self-criticism stage his voice is almost swallowed by ``an attitude problem''.

He has at least three reasons. First, saving an out-of-bounds ball leaves no time for hesitation; half a second loses it. Players know such a sprint contains no separate procedure for checking nearby people. Demanding observation, decision, and braking within three seconds effectively demands that he not save the ball. Second, he accepted the consequences: helping the boy up, gathering books, and apologising immediately. The boy said it was fine. The later assignment demands that an accident become wrongdoing, which is rewriting facts, not reflection. Third, and most painfully, five accounts circulate while he lacks a channel. The younger boy has infirmary evidence, spectators have phones, and the representative has authority to report to the teacher. Xiaoming's version was condemned when ``self-criticism'' was assigned, because the sorrowful self he must perform and the self saying ``I was not wrong, but I am sorry'' are different people.

These arguments are substantial. You may reject his conclusion, but cannot simply say he refuses to admit anything. What he accepts and what the assignment requires him to accept are different.

Pause and count the voices: injured boy, collider, spectator, reporter, sharer. Complete? Someone is missing.

\textbf{The person everyone discusses but nobody seriously consults is the original seventh-grade boy.} Did he not speak? Yes, but only one sentence: ``He came straight at me.'' Afterwards, everyone borrows his knee and glasses as evidence. The teacher uses them to show Xiaoming wrong; his friends to show it was not serious; the chat group to show older pupils bullying younger ones. His injury is everywhere; his voice disappears. He is an \textbf{observed person}, watched and narrated by everyone. What he might want to say, such as wanting to watch the game and feeling warmed when someone helped him up afterwards, enters no version.

Scale this up to history and it shapes almost all the past we can read.

In West African Mande societies, \textbf{griots} are hereditary custodians of memory. From childhood they learn centuries of genealogies, battles, and oaths, singing them with instrumental accompaniment at ceremonies. Consider the implication. Without written records, the past resides in particular people, who are court employees with positions of their own. A community's historical sound depends largely on whom its griot serves. Royals with griots possess their version of history; those without lack even the privilege of being misrepresented and remain silent.

Turn to the Americas. Enslaved Black people in the nineteenth century were almost entirely denied narrative authority. Accounts of slavery long belonged to enslavers and their sympathisers. After escaping, Frederick Douglass wrote his experiences himself in 1845. His autobiography shook America. One detail may surprise you: two white abolitionist leaders supplied prefatory endorsements testifying that this Black man really wrote the book. A participant telling his own story required others to guarantee his ability to tell it. Such inequality meant he had to prove he deserved to speak before speaking.

The most extreme dramatisation of narrative authority as life itself is the frame of \emph{The Thousand and One Nights}. Betrayed by his wife, Shahryar marries a bride daily and kills her at dawn. Scheherazade, the vizier's daughter, volunteers to enter his palace and tells stories nightly, reaching their gripping points at sunrise. To hear more, he lets her live. She tells for three years and lives for three years. Storytelling is not a talent show but her sole means of survival. Each night answers who speaks: someone initially denied even the right to survive until dawn. Speaking changes who deserves to live.

Hearing multiple voices is therefore more than asking everyone present in turn. It means noticing whose voices are placed outside the story and who remains merely a prop in others' accounts. A class, a community, and an age all face this problem.

\section[Thinking Bridge: Three Narrative Questions]{Thinking Bridge: Who Speaks, Who Sees, Who Is Silent?}
Can we analyse a self-criticism, news report, history, chat screenshot, or film with a portable tool instead of intuition alone? Yes. Narratology's core tool is three questions asked in sequence:

\textbf{First: who is speaking?}

Separate author and narrator. Who wrote the words, and whose voice speaks inside them? In which grammatical person? Beware of devices making you forget the question. Self-criticism wants its sorrowful paper self mistaken for real Xiaoming. News releases saying ``it is understood'' or ``a relevant official stated'' discourage asking whose mouth stands behind the phrase. Separate author, narrator, and character, and many claims of sincerity or objectivity reveal seams.

\textbf{Second: who is seeing?}

Speaking and seeing can separate. The narrator speaks, but whose eyes supply the scene? This is \textbf{viewpoint}. Lock the court story to the younger boy's eyes, and nothing warns of the flying ball; Xiaoming initially appears only as a tall figure. Lock it to Xiaoming's eyes, and the reader sees the escaping ball first, discovering someone at the sideline only in peripheral vision during the sprint. One collision distributes surprise differently. The self-criticism is stranger still. It pretends to use my eyes while actually borrowing the teacher's. Xiaoming must look back down at himself from the office to write ``disgraced our class''. Being required to see oneself through another's eyes is the genre's real training.

Incidentally, omniscience, knowing everything and entering every mind, does not automatically mean objectivity. It hides selection more deeply. A camera claiming access to every mind still chooses whom to inspect first, later, or never. We will return to this easily misjudged point.

\textbf{Third: who has not spoken?}

This is the sharpest question. Who is described, watched, and classified without receiving a sentence of their own? The younger boy in the court story; all actual participants in the chat's bullying version; often defeated people, women, servants, and children in histories. The more complete and fluent an account, the more carefully you should look for silences trimmed away.

Try all three on any school reading. Who is the author, and how do they relate to the text's ``I''? Whose eyes guide the camera? If briefly mentioned characters could speak, how would the story change? Many familiar passages suddenly grow another story. That is the tool's power.

\section{A Madman's Written Truth}
Read a short primary passage. In 1918, Lu Xun published \emph{Diary of a Madman}, modern Chinese literature's first vernacular story, presented as a persecutory-delusion sufferer's diary. Its best-known passage reads:

\begin{quote}
I opened history to examine it. This history had no dates; across every crooked page were the words ``benevolence, righteousness, and morality''. Unable to sleep, I read carefully for half the night, until I found words between the lines. The entire book said ``eat people''! (Lu Xun, \emph{Diary of a Madman})
\end{quote}
Apply the questions. Who speaks? The narrator is the madman, convinced everyone wants to eat him; the author is Lu Xun, a clear-headed writer. Their distance powers the story. In 1918, directly accusing China's old ethics of eating people would bring serious trouble. He therefore creates a madman's mouth. Every accusation literally comes from a patient's ravings and can be dismissed as madness.

Another mechanism is often overlooked. Before the diary stands a short classical-Chinese preface in an outsider's voice, explaining that the diarist recovered and went elsewhere to await an official appointment. Once cured, he seeks office. Thus the words exposing morality's cannibalism are classified as illness, while recovery means rejoining the system he saw through while ill. Lu Xun quietly exchanges madness and normality.

Here is an important counterexample. We assume reliable narrators tell truths and unreliable ones falsehoods. The madman reverses this: \textbf{the work's only unreliable narrator is also its only truth-teller}. His apparently normal, impeccably spoken brother, doctor, and neighbours are all reliable, all protecting the cannibalistic order. Unreliable does not mean liar. A disordered voice can see through everything; a respectable voice can build a prison with every sentence. Reliability is judged not by fluency or a teller's apparent normality but by factual testing and whose interests the account serves.

\section{Three Confessions in a Grove}
What if several accounts of one event contradict one another so thoroughly that they cannot all be true? This is no puzzle game; it happens daily in courts, news, and families. Ryūnosuke Akutagawa's 1922 story ``In a Grove'' pushes it to the limit. Here is its outline.

A samurai's body is found in a grove. Investigators hear several statements. The bandit Tajōmaru admits killing him: he subdued the samurai, then defeated him honourably in a fair duel. The wife admits killing him: after her humiliation, her husband looked at her with cold contempt, and in confusion she stabbed him with a dagger. Strangest of all, the dead samurai speaks through a medium and claims suicide. His wife yielded to the bandit and asked him to kill her husband, so the samurai despairingly took his own life. All three compete to be the killer. Each account is detailed, emotional, and places its teller in a barely coherent moral position: honourable duellist, pitiable woman destroyed by a look, or proud betrayed husband.

Factually we hold only a body, a blade, and a grove. Here are three genuinely different explanations for the clashing confessions.

\textbf{Explanation one: memory itself is unreliable.}

Psychological experiments repeatedly show memory is reconstruction, not video replay. Emotion, position, and later information enter every recollection. Witnesses describe one accident in radically different ways while each sincerely believes their memory accurate. Here the contradictions are defects in cerebral hardware: nobody deserves blame; nobody remembers perfectly. The account is generous and experimentally supported. Its limit is that it cannot explain why every memory's error happens to beautify its narrator. Random errors should not align so consistently.

\textbf{Explanation two: everyone is defending themselves, not misremembering.}

The confessions become three emergency operations to rescue self-image. Memory is raw material; interests are scissors. This sharp account explains the direction of errors but makes people too calculating. It assumes each knows the truth and deliberately conceals it. Yet people often sincerely embellish themselves and genuinely feel wronged. Calling everything performance cannot explain sincere error.

\textbf{Explanation three: truth cannot be reached outside telling.}

More radically, we possess only narratives. Once facts pass through any mouth, they acquire its grammatical person, viewpoint, and editing. No ownerless account exists. Asking who provides the original may be mistaken; ask instead what each narrative accomplishes. This explains why even honest witnesses cannot offer neutral stories. But its danger is greatest. If everything is narrative, where do hard facts such as the samurai's death and the blade in his chest belong? People who call everything narrative often seek to evade hard facts. Keep the boundary clear: narratology says telling always has a viewpoint, never that facts do not exist.

Each account grasps part of the truth, none all of it. Akutagawa's most ruthless choice is to keep an objective, omniscient narrator absent throughout. No voice tells you what really happened. You finish with voices alone, judgement returned to your hands.

\section[Further Challenge: Separate Voice and Vision]{Further Challenge: Separating Speaking from Seeing}
Now bring in the chapter's weightiest theory. In the 1960s and 1970s, French narratologist Gérard Genette made a seemingly simple but profoundly influential distinction: \textbf{who speaks, narrative voice, and who sees, focalisation}.

Earlier terminology combined both as viewpoint. Genette treated them as separate accounts. Speaking concerns voice; seeing concerns the presented scene. A third-person narrator outside the story can stay close to a character's eyes inside it, as in ``Xiaoming stood outside the office. Now he was done for.'' The narrator speaks; Xiaoming sees. Conversely, a first-person narrator can deliberately conceal something their earlier self saw. Voice and scene are independent knobs, one in each storyteller's hand, producing countless effects.

Genette distinguishes three forms of focalisation by whether the narrator knows more or less than characters:

\textbf{Zero focalisation}: the narrator knows everything, everyone's thoughts, past, and future, commonly called omniscience. Much classical Chinese fiction works this way, letting storytellers enter any mind. Its strength is freedom; its weakness is the counterexample planted earlier. Omniscience resembles God's viewpoint, but a narrator must choose whom to inspect first. \textbf{Omniscience is not objectivity; it hides selection within the composure of knowing everything.} Mistaking omniscient third-person telling for closeness to fact is a common reading error.

\textbf{Internal focalisation}: the scene stays within one character's senses; readers know only what that person knows. Detective fiction commonly does this. You follow the detective, cannot see what they cannot see, and know less than the murderer until the final chapter. Self-criticism is a twisted instance: nominally the camera stays inside my eyes, but orders require it to move to the teacher's position.

\textbf{External focalisation}: the narrator knows less than characters, recording outward speech and action without entering minds, like an unfeeling roadside camera. ``He opened the door, sat down, read the letter three times, then burned it.'' What was he thinking? The narrator will not say. Interpretation falls entirely to the reader.

A companion tool comes from Wayne Booth's 1961 \emph{The Rhetoric of Fiction}: the \textbf{implied author}. Writing creates a second self. Real Xiaoming is messy, lazy at times, and capable of lying, but once he writes, even a self-criticism, the total choices form an image of Xiaoming on paper. Readers know that implied author, never the actual person. Hence reading a writer's complete works can make you feel you know them, only to disappoint you when meeting them. You knew the implied author. The teacher likes the properly contrite implied Xiaoming; Xiaoming resents being responsible for an image that is not himself.

Booth's other tool has already appeared in the madman and grove: the \textbf{unreliable narrator}. When a narrator conflicts with other textual evidence, readers are invited to doubt the voice. This sophisticated device lets the author lie on one level so readers see truth on another.

The device is older than you may think. In the \emph{Odyssey}, Odysseus returns after ten years of wandering, and the four most thrilling books, the Cyclops, sorceress, and Sirens, are entrusted to his own telling at a palace banquet. The most marvellous adventures therefore have no omniscient narrator's guarantee. They come from a famously cunning man's self-account, heard by strangers newly moved by his stories. Audiences over two thousand years ago occupied our seat: the more movingly he speaks, the more we must ask whether he has witnesses. The epic never answers, leaving the question inside the song.

Now the title's full weight is clear. Who speaks concerns voice; who sees concerns the scene. Skilled narrative manipulates their gap to decide where sympathy falls. \textbf{Distributing viewpoints distributes sympathy.} Readers' hearts follow the camera. Whether a court hears a victim first or a defendant's explanation, a film gives a villain three minutes of childhood memories, or a textbook begins a war with one side's campaign diary all quietly allocate moral judgement. Narratology is self-defence for every listener, not merely a literary enthusiast's craft.

\section{After Everyone Gets a Microphone}
These ancient problems explode daily inside your phone, wearing new clothes.

First, apparent good news: technology seems to solve silence. The seventh-grader has a phone, campus posting board, and group chat. Douglass, who two centuries ago needed white endorsements for publication, could now simply register an account. Everyone has a microphone and can become narrator.

But new difficulties follow.

\textbf{The grove in the age of lengthy personal posts.} Parties to a trending dispute publish successive long statements, each richly detailed, emotional, and tightly self-defending. Commenters swing between them. You recognise the structure: ``In a Grove'', with seven testimonies becoming seven million shares. The online advice to wait before judging means, in narratological terms, refusing to treat one internally focalised account as omniscience before all accounts arrive.

\textbf{Screenshots and editing power.} Chat screenshots are among today's commonest evidence, but inherently have beginnings and endings cut off. Whoever captures them gains editing authority. Apply the third question: who has not spoken? The party whose context was removed.

\textbf{People being watched.} Cameras, street photographs, and casually uploaded images create new observed people. A courier crying in a lift receives an inspirational caption; a stranger's sleeping posture on a train becomes online ridicule. They did nothing wrong, merely entered someone else's lens. They inhabit countless stories without narrative authority over any. Viewpoints distribute sympathy, and cameras are never allocated randomly.

\textbf{The hardest problem comes next.} Your conversational partner, or a first-person account you read, may not be human-written at all. Machines mass-produce heartfelt selves. The self-criticism problem, the paper's I not being a real person, becomes everyday internet life. The once-easiest question, who speaks, becomes the hardest. Precisely for that reason, all three deserve more practice than ever.

\section[For You: Which Account Enters the Archive?]{For You: Which Account Would You Archive?}
Return to the court.

Imagine the school creates an archive of the collision, sealed for fifty years. Only one account may become its sole official record. Before you lie five documents: the boy's infirmary statement, Xiaoming's self-criticism, the girl's recollection, the sports representative's report, and screenshots of the year-group chat.

Which do you preserve, and why?

Before deciding, review your tools.

You might choose the self-criticism, but its paper self stands behind a performance. The boy's statement, perhaps, but its internal viewpoint misses the out-of-bounds ball. The bystander's account, perhaps, but seeing most and lacking a position differ. You might compromise and include all five, but the rule permits one. Choosing allocates sympathy for fifty years of future readers.

A harder question follows. If this record determines how people react to Xiaoming's name, would you ask instead: \textbf{who chooses the document, and should that power of selection be recorded before the document itself?}

Take it personally. If you become a prop in someone else's story, narrated, classified, or screenshotted, what rules should the storyteller follow? Write them down. Then ask whether you follow them when telling others' stories.

There is no standard answer. But from today, whenever you hear a story, look again at its teller: where they stand, where the camera points, and whose silence has been trimmed away.
